\chapter{Colorectal Cancer Stool Screening}

\section*{What Is This Test?}
Stool-based colorectal-cancer screening tests look for hidden blood, abnormal DNA, or both in a stool sample. The main noninvasive options are the fecal immunochemical test (FIT), guaiac fecal occult blood test (gFOBT), and multitarget stool DNA-FIT test. These tests screen for colorectal cancer and some precancerous lesions; they do not diagnose the cause of an abnormal result.

\section*{Alternative Names}
Fecal Immunochemical Test; FIT; iFOBT; Fecal Occult Blood Test; FOBT; gFOBT; Stool DNA Test; FIT-DNA; Multitarget Stool DNA Test; Cologuard

\section*{Why It Is Ordered}
Average-risk colorectal-cancer screening according to age and local recommendations, or selected follow-up when a clinician chooses a stool-based strategy. The appropriate interval depends on the test: FIT and gFOBT are generally repeated more often than stool DNA-FIT. People with symptoms, a personal or strong family history of colorectal cancer or polyps, inflammatory bowel disease, or hereditary cancer risk may need a different evaluation.

\section*{How This Test Is Performed}
The patient collects stool at home using the kit supplied by the laboratory or provider and returns or mails the sealed specimen according to the instructions. FIT detects human hemoglobin. gFOBT uses a chemical reaction to detect heme. Stool DNA-FIT combines molecular testing for selected DNA changes with an immunochemical blood test.

\section*{Preparation, Risks, and Limitations}
FIT usually requires no dietary restriction, while gFOBT may require medication or dietary instructions. Follow the kit instructions precisely and avoid mixing the specimen with urine or toilet water. The tests are noninvasive and have no procedure-related risk. A positive result does not prove cancer, but it generally requires diagnostic colon evaluation arranged by a clinician. A negative result does not exclude all polyps or cancers, so screening must be repeated at the recommended interval.

\section*{References}
\begin{enumerate}
  \item MedlinePlus. Colorectal Cancer Screening Tests.
  \item United States Preventive Services Task Force. Colorectal Cancer: Screening.
\end{enumerate}

\clearpage
